\documentclass[11pt]{article}
\usepackage[margin=0.95in]{geometry}
\usepackage{amsmath,amssymb}
\usepackage{graphicx}
\usepackage{booktabs,longtable,array}
\usepackage{xcolor}
\usepackage[hidelinks]{hyperref}
\usepackage{enumitem}
\usepackage{caption}
\setlist{itemsep=2pt,topsep=4pt}
\captionsetup{font=small,labelfont=bf}
\setlength{\parskip}{5pt}
\setlength{\parindent}{0pt}
\renewcommand{\arraystretch}{1.15}
\title{Phase 2 Summary: Dimensionality, Coupling and the Re-analysis}
\author{}
\date{trm-delay-robust-scaling $\cdot$ phase summary $\cdot$ 2026-09-09/10 $\cdot$ 54 runs\\[2pt]Labels are defined in \texttt{.\allowbreak{}.\allowbreak{}/\allowbreak{}glossary.\allowbreak{}md}.}
\begin{document}
\maketitle
\section*{Question}



Does a different kind of difficulty, high-dimensional coupled dynamics, reward recursion? And is any null about recursion in general or about this particular recursive structure?



\section*{Design}



\begin{itemize}
\item \textbf{Environment:} BipedalWalker-v3 (24-D observation, 4-D action, coupled multi-joint dynamics).
\item \textbf{Trunks (6):} MLP, Transformer, LSTM, TRM, plus two TRM variants requested after Phase 1 review:
\begin{itemize}
\item \texttt{trm\_hier6}, with six parallel low-level latents;
\item \texttt{trm\_singlez}, with a single latent.
\end{itemize}
\item \textbf{Conditions (3):}
\begin{itemize}
\item no delay;
\item fixed d = 1 + queue;
\item fixed d = 1 without the queue (the augmentation ablation, C2).
\end{itemize}
\item \textbf{Seeds:} 3, so 54 runs on V100 GPUs. A filename collision silently overwrote three MLP cells; it was caught before analysis and six jobs were re-run.
\end{itemize}



\section*{Results}



\begin{itemize}
\item \textbf{No comparison reached significance in any direction:} TRM $-$ MLP 2/3, TRM $-$ Transformer 0/3, TRM $-$ LSTM 1/3.
\item \textbf{The recursive structure made no detectable difference.} Base, hier6 and singlez reversed order from seed to seed.
\item \textbf{Augmentation had no clear effect.}
\item \textbf{Cost:} TRM 5.2$\times$, hier6 6.9$\times$, singlez 8.0$\times$ the MLP's wall-clock.
\end{itemize}



\section*{The post-hoc re-analysis (all 102 Phase 1--2 runs, no new compute)}



It changed what both phases meant.



\begin{enumerate}
\item \textbf{Two of three environments were floor effects.} MountainCar was exploration-limited (median $-$8.9 against a solved +90, the documented DDPG + OU-noise failure). No BipedalWalker run learned to walk (best $-$81.1 against +300). \textbf{The complexity question had never been put at risk: untested, not disconfirmed.}
\item \textbf{TRM fit the value function worse} (critic loss higher in 7/7 cells early, median 2.72$\times$, p = 0.008) at unchanged learning speed (117 vs 115 episodes to 90\% of span). This was explanation E1, an optimisation/conditioning deficit.
\item \textbf{The deficit looked like capacity.} The 797k Transformer showed it too and the 167k LSTM did not.
\item \textbf{Three confounds}, any one of which could explain the Phase 1 negative:
\begin{itemize}
\item \textbf{C-1, capacity:} 788k vs 135k parameters.
\item \textbf{C-2, latent reset:} the port re-initialised its latent every call, HRM-Agent's weaker variant.
\item \textbf{C-3, nothing to infer:} under fixed delay + queue the observation is already sufficient.
\end{itemize}
\item \textbf{Preconditions P1--P3} for TRM's in-domain success, which the port lacks: refinement continuity, a correctable structured answer, per-instance capacity.
\end{enumerate}



The report was rewritten after author review rejected the first draft, then the literature review was integrated (26 sources). That retracted the novelty claim, because HRM-Agent precedes this work.



\section*{What it established}



The honest position after Phase 2 was not "recursion does not work". It was "the experiment that would tell us has not been run yet". Removing C-1, C-2 and C-3 became Phase 3's job. The recommended experiments E1--E5 became Phase 3 Stages 1--2 and Phase 4 Stage A.
\end{document}
